\documentclass{article}
\usepackage{amsmath,amssymb,hyperref}
\title{Shared interaction derived from native hom products}
\date{15 November 2026}
\begin{document}\maketitle
\paragraph{Status.} General construction independently reviewed under the stated hypotheses. This derives a same-state joint-witness operation from categorical products, rather than assuming its desired whole-network law. It is distinct from independent-state tensor parallel and from nondeterministic choice. The distinction of their units is essential.

\paragraph{Independent data.} Let $\mathcal K$ be the process bicategory and $\mathcal S$ the retained state bicategory, with the original pure-action pseudofunctors $E_A:\mathcal S\to\mathcal K$. Assume each needed hom category of $\mathcal K$ has specified finite products. For every admitted state map $h:X\to Y$, require pure input precomposition
\[
 (-)E_Ah:\mathcal K(E_AY,E_BY)\to\mathcal K(E_AX,E_BY)
\]
to preserve these products through its canonical comparison. This is a condition on the pure interface before joint machines are defined. The component state map need not be invertible. No finite-product preservation of output postcomposition $E_Bh(-)$ is assumed.

For any functor $V$ between categories with products there is a canonical comparison $V(\prod_i f_i)\to\prod_i Vf_i$, defined by its projected maps $V\pi_i$. It is natural, unital and compatible with composition of functors by product uniqueness. These facts, not chosen arbitrary comparison cells, supply the forward part of the construction.

\paragraph{Common-state joint operation.} Let $M_m(A,B)$ have one state $X$, bodies $f_i:E_AX\to E_BX$, one state map $h$ with a native square $\alpha_i$ for every $i$, and one state two-cell satisfying all native equations. This is the explicit common-state multi-operation doctrine, not a product of independently allocated machines. For finite $m$ define
\[
 \mathsf J_m(X,(f_i))=(X,\prod_i f_i).
\tag{1}
\]
Its native square on $(h,(\alpha_i))$ is
\[
 E_Bh(\prod_i f_i)
 \longrightarrow\prod_i(E_Bh f_i)
 \xrightarrow{\prod_i\alpha_i}\prod_i(g_i E_Ah)
 \xrightarrow{\cong}(\prod_i g_i)E_Ah.
\tag{2}
\]
The first arrow is canonical and may be noninvertible. The last is the inverse canonical comparison justified by the independent input-preservation hypothesis. No general process interchanger is inverted.

Keep the exact $h$ and the exact state two-cell on native cells. To check its native equation, postcompose both routes with every product projection, using the invertible final comparison in (2). They reduce to the original component equations and naturality of the canonical product comparison. The same test proves preservation of composition and identities, with the actual $E_A,E_B$ pseudofunctor frames inserted. Thus (1) is the proposed native pseudofunctor over the identity state base. Product associativity, units and finite substitution comparisons lift over identity state maps by their canonical process isomorphisms. Their naturality and equations are detected by product projections, not by a potentially nonfaithful pure realization.

\paragraph{A derived sequential interchanger.} For four bodies on one state, $f_i:E_AX\to E_BX$ and $g_i:E_BX\to E_CX$, the two pairs of product projections give cells
\[
 (g_1\times g_2)(f_1\times f_2)\longrightarrow g_i f_i.
\]
Universality gives the forward comparison
\[
 (g_1\times g_2)(f_1\times f_2)
 \longrightarrow(g_1f_1)\times(g_2f_2).
\tag{3}
\]
This is an actual native arrow over the identity state map between the corresponding closed machines. It is not assumed invertible. For a common state map and four native squares, its naturality follows by projecting to each target component and using the original square-pasting equation. For multiple chronological layers, every expanded route projects to the same chronological composite of the chosen projections. Product uniqueness and process associativity prove the composition and unit equations for this directed comparison. Finite arities use the same diagonal selection of a common component index.

The nullary joint process is the terminal object of the relevant process hom category. It is generally not the chronological identity $1_{E_AX}$. Thus this structure must not be called the original normal tensor parallel with unchanged units. It is a separate local, conjunctive or joint-witness interaction, with the displayed directed composition comparison. Its nullary comparison is induced by the terminal process universal property. No unproved duoidal or concurrent-pseudomonad axiom inventory is inferred from that description.

\paragraph{Synchronization: directed in general.} Assume native cones in the common-state multibody doctrine and output dialgebra are created over the same selected state cones; in particular the relevant $E_B$ preserves them. Apply $\mathsf J_m$ to the universal multibody cone. This produces an actual native lax cone, since (1) is a pseudofunctor. Output universality then gives
\[
 \mathsf J_m L_J^{M_m}F\longrightarrow
 L_J^{D(A,B)}(\mathsf J_mF).
\tag{4}
\]
Its coherence follows from the representing cone equivalence and the already admitted operation comparisons. It is generally only directed: (2) may fail to preserve pseudo projection squares because its output comparison may be noninvertible. The identity state base alone does not justify using the strong-preservation criterion.

If, additionally, the output postcomposition functors for the actual chosen cone projection state maps preserve the finite products at issue, then (2) takes those pseudo projections to pseudo projections. The existing creation/preservation argument recognizes the transported cone as universal and makes (4) an equivalence. This is a sufficient exactness condition on those projections, not an assertion about arbitrary state maps or all physical systems.

\paragraph{A native non-thin source.} Let $\mathcal C$ have finite limits and take the process bicategory $\operatorname{Span}(\mathcal C)$ with all apex maps as native cells. A hom category is $\mathcal C/(A\times B)$. Its product is the pullback of witness apices over $A\times B$; its terminal object is the universal span with apex $A\times B$. Hence a joint transition retains a tuple of actual compatible witnesses with common input and output. It is not an isomorphism-class or relation quotient, and repeated factors need not be idempotent.

For graph-pure input maps, precomposition is base change on these slices, which preserves finite products by pullback pasting. Output graph composition is direct image by composition of the slice structure map; it has the canonical forward product comparison but need not preserve products. For example in finite sets, two spans from $1$ to $\{0,1\}$ supported respectively at $0$ and $1$ have empty product. After collapsing the output set to $1$, their images have a singleton product. The output comparison is the noninvertible map from the empty witness set to the singleton. This verifies both the source hypotheses and why the strong synchronization claim requires extra exactness. A state action by product with external interfaces uses these same graph-pure maps and retains the original tags.

\paragraph{Established source and exact use.} Lack, Walters and Wood, \emph{Bicategories of spans as cartesian bicategories} (2010), introduction and Section 2, \url{https://arxiv.org/pdf/0910.2996}, explicitly identify span homs with slice categories and distinguish their local products from the external tensor. These structures are prior art, not a newly discovered parallel formalism. The present calculation integrates the local product with the given native state tags and the separately constructed synchronization doctrine. It does not infer that every relative factory is a cartesian bicategory or has all adjoints of pure maps.

\paragraph{Meaning and limits.} The source derives a particular compatible-witness conjunction. It does not force arbitrary processes to share a realizable physical transition, choose a stochastic coupling, or enumerate schedules. The product may have no witnesses. Chronological same-state composition forwards updates; local product instead requires compatible endpoints. Their interaction is exactly the forward map (3), which can forget the requirement of one common intermediate witness state. This explains its direction without assuming independent effects commute.

The construction supplies a general source-derived same-state interaction law and a native directed synchronization comparison. It does not identify this law with the original tensor, change its units silently, solve arbitrary instantaneous feedback, or supply missing product-preservation hypotheses in the original factory. All further feedback equations must use the actual admitted joint-witness operation and its native data.
\end{document}
